\documentclass[a4paper,12pt]{article}
\usepackage{bbding,combelow,textcomp,amsfonts,amsthm,amsmath,amssymb,graphicx,color,hyperref,etoolbox}
\usepackage[utf8x]{inputenc}
\usepackage[lined,boxed,commentsnumbered]{algorithm2e}

\newtoggle{story}


%%%%%%%%%%%%%%%%%%%%%%%%%%%%%%%%%%%%%%%%%%
%% Customisation options --- update as necessary
%%%%%%%%%%%%%%%%%%%%%%%%%%%%%%%%%%%%%%%%%%

%% Course details: title, semester, instructors

\newcommand{\courseTitle}{Introduction to Probability Theory (Math 2502)}
\newcommand{\semester}{Spring 2026}
\newcommand{\instructorA}{Shagnik Das}
\newcommand{\instructorB}{}

%% Sheet number

\newcommand{\sheetNumber}{3}


%% Submission details

\newcommand{\dueDate}{13:00, Apr 7th.}
\newcommand{\lateWarning}{Submissions that are late will receive no credit.}


%% Display irrelevant footnotes?  \toggletrue for yes, \togglefalse for no

\toggletrue{story}

%%%%%%%%%%%%%%%%%%%%%%%%%%%%%%%%%%%%%%%%%%
%% End of customisation options
%%%%%%%%%%%%%%%%%%%%%%%%%%%%%%%%%%%%%%%%%%

\pdfpagewidth 8.5in
\pdfpageheight 11in
\topmargin -1in
\headheight 0in
\headsep 0in
\textheight 8.5in
\textwidth 6.5in
\oddsidemargin 0in
\evensidemargin 0in
\headheight 77pt
\headsep 0in
\footskip .75in

\makeatletter
\newcommand{\buildtitle}[4]{
\begin{flushleft}
{\large
#1
\hfill{}
#2
\par
#3
}
\end{flushleft}
\vskip 4pt
\begin{center}
{\large\bfseries#4\par}
\end{center}
\bigskip
}
\makeatother

\renewcommand{\thesection}{\normalsize\arabic{section}}

\renewcommand{\deg}{\textrm{deg}}
\newcommand{\eps}{\varepsilon}
\newcommand{\ceil}[1]{\left \lceil #1 \right \rceil}

\newcommand{\hint}{\begin{flushright} [Hint at \url{\hintURL}.] \end{flushright}}

\newcommand{\card}[1]{\left| #1 \right|}
\newcommand{\mb}[1]{\mathbb{#1}}
\newcommand{\mc}[1]{\mathcal{#1}}


\newcommand{\story}[1]{\iftoggle{story}{\footnote{#1}}{}}
\newcommand{\storymark}[1][42]{\iftoggle{story}{\footnotemark[#1]}{}}
\newcommand{\storytext}[2][42]{\iftoggle{story}{\footnotetext[#1]{#2}}{}}

\newcommand{\bonus}[2]{\paragraph{Bonus (#1 pt\ifstrequal{#1}{1}{}{s})} #2}


\begin{document}


\buildtitle{\courseTitle{}}{\semester}{\instructorA{} \hfill \instructorB{}}{Exercise Sheet \sheetNumber{}}

\vspace{-0.2in}

\begin{center}

{\bf Due date: \dueDate \\ \lateWarning}

\end{center}

\noindent You should try to solve all of the exercises below --- each problem is worth $10$ points. Make sure to clearly justify your answers, defining everything precisely and stating any results you are using. You should then submit your solutions in Gradescope on NTU COOL, either as a typed PDF or a clear scan of your handwritten answers, and should indicate where in the file the solutions to each exercise are. If you have difficulties with Gradescope, please send your solutions by e-mail.

\paragraph{Exercise 1}  Let $X$ be a random variable that takes nonnegative integer values. Prove that
\[ \mb E [X] = \sum_{n \in \mathbb{N}} \mb P (X \ge n). \]

\paragraph{Exercise 2}  A store is offering a promotion --- they have $n$ different types of coupons, and every time you buy something, you get a coupon of a uniformly random type, independently of your previous coupons. Once you have collected all the different types of coupons, you can redeem them for a free probability textbook!

Recall that in lecture we saw that the probability you will need to collect more than $t$ coupons before you have a full set is
\[ \sum_{i = 1}^n (-1)^{i+1} \binom{n}{i} \left( \frac{n-i}{n} \right)^t. \]

What is the probability that, after collecting $t$ coupons, you will have exactly $k$ different types in your collection?

\paragraph{Exercise 3}  A casino introduces a new game for $n$ players. The dealer has a deck of $n+1$ cards, labelled from $1$ to $n+1$, which she shuffles. She then gives each player a (different) random card, leaving the remaining card for herself.

The dealer will then reveal her card, after which the players reveal their cards, one at a time. Every player whose card is higher than all of the previously revealed cards wins \$100.

How much should the casino expect to pay out in prizes each round?

\paragraph{Exercise 4}  An absent-minded professor has $10$ keys on his keyring, labelled $1$, $2$, ..., $10$ in cyclic order. Key $6$ opens his front door, but he has unfortunately forgotten this crucial piece of information.

He starts by trying Key $1$, which of course does not open the door. He then tries a neighbouring key --- either the next or the previous one --- uniformly at random (in this instance, the second key he tries could be Key $2$ or Key $10$).

The professor is unable to remember which keys he has already tried, so every time a key doesn't work, he next tries one of its two neighbouring keys, uniformly at random, independently of all previous choices.

What is the probability that he tries all the other keys before he tries Key $6$ and is finally able to open the door?

\newpage

\section*{Hints} % The following QR codes contain hints for some of the homework exercises.  You should be able to decode them using any QR code scanner, including this one: \url{https://zxing.org/w/decode.jspx}.

\paragraph{Exercise 1} Also available at \url{https://i.ibb.co/GywL4X4/S03E1.png}.

\begin{center}
	\includegraphics[scale=0.5]{Hints/S03E1.png}
\end{center}

\paragraph{Exercise 3} Also available at \url{https://i.ibb.co/B249SB5Z/S03E3.png}.

\begin{center}
	\includegraphics[scale=0.5]{Hints/S03E3.png}
\end{center}

\paragraph{Exercise 4} Also available at \url{https://i.ibb.co/yB0GQqbd/S3E04.png}.

\begin{center}
	\includegraphics[scale=0.5]{Hints/S03E4.png}
\end{center}

\end{document}